% -----------------------------------------------------------------------------
% Discussion
% -----------------------------------------------------------------------------
\section{Discussion}
\label{sec:discussion}
The three research questions posed in Section~\ref{sec:introduction} can now
be answered directly. For RQ1, under matched input content the sequence
models (ODST and attention--linear) matched or exceeded the flattened
references. The order-perturbation and partial-history experiments
showed that the frozen sequence models rely on both chronological order and
accumulated behavioural history rather than input volume alone. For RQ2,
teacher-anchored logit-and-route regularisation produced reproducible
end-to-end student optimisation across the three predefined r5.2 seeds, with
students closely matching their frozen teachers and limited prediction and
routing drift. For RQ3, the architecture provided faithful reference-centred
tree-contribution evidence, improved probability calibration through grouped
temperature scaling, an interpretable alert-consolidation view, a
component-level latency profile identifying the ODST head as the runtime
target, and reproducible degraded-input sensitivity references. The
following subsections examine each in turn.

\subsection{Temporal Representation and Predictive Position}
\label{subsec:discussion_temporal}
The chronological representation preserved inactive days and prevented
sequence windows from crossing partition boundaries. It also retained the order
and duration of behavioural activity. The order-perturbation and
partial-history experiments confirmed that both chronological arrangement
and accumulated context contributed to prediction. Performance increased
as the available history expanded from one to 20 days, while reversing or
shuffling the same daily observations reduced discrimination. The value of
the representation therefore arises from temporal structure rather than
from window size alone.

\smallskip
The same-information comparison further separated temporal representation
from input quantity. The oblivious differentiable sparsemax tree (ODST) and
attention--linear models both achieved a mean average precision (AP),
reported as PR-AUC, of 0.931, with F1-scores of 0.910 and 0.908. Sequence
models also achieved metric-specific advantages over several flattened
references. These results position the bidirectional long short-term memory
(Bi-LSTM)--attention pathway as the principal source of temporal predictive
value.

\smallskip
The multi-day cohort provided complementary evidence for sequences
containing verified malicious activity on more than one day. ODST remained
competitive across the evaluated models, with a supported cohort PR-AUC
advantage over Random Forest (RF). This result strengthens the case for
preserving distributed behavioural context without converting the cohort
into a universal definition of low-and-slow activity.

\smallskip
Strong RF and Extreme Gradient Boosting (XGBoost) results are scientifically
expected when behavioural records have already been converted into
structured tabular inputs. Grinsztajn, Oyallon, and Varoquaux showed that
tree ensembles remain demanding references on typical tabular data
\cite{Grinsztajn2022}. Their benchmark used balanced classification and
excluded stream-like and time-series tasks, so it provides useful context
for the baselines rather than a direct prediction of CERT performance.

\subsection{Teacher-Anchored End-to-End Optimisation}
\label{subsec:discussion_teacher_anchoring}
The architectural development progressively identified routing stability as
the central optimisation requirement. The original differentiable forest
established a shared gradient pathway through the temporal encoder,
attention mechanism, and tree classifier. Subsequent diagnostics showed that
unrestricted joint updates could preserve useful temporal representations
while concentrating ODST routes and leaves.

\smallskip
Teacher anchoring addressed this requirement by using stable prediction and
routing references during student optimisation. The complete
Bi-LSTM--attention--ODST student remained trainable, while the teacher was
removed after training. The locked procedure reproduced across CERT r5.2
seeds 42, 52, and 62. Student PR-AUC values closely matched their frozen
teachers, with a mean student-minus-teacher difference of approximately
\(+0.00043\), interpreted as predictive non-degradation rather than
improvement.

\smallskip
The importance of this result lies in reproducibility and end-to-end
viability across the three predefined r5.2 seeds, not in cross-version
transfer: teacher-anchored training was performed only on r5.2, and no
equivalent teacher-anchored procedure was run on r4.2. The student preserved
the functional behaviour of the frozen model while permitting genuine
updates to the encoder, attention mechanism, and ODST head. Teacher
anchoring therefore provides the optimisation contribution of the study
rather than a mechanism for producing a large score increase.

\subsection{Structured Explanation, Calibration, and Alert Interpretation}
\label{subsec:discussion_explanations}
The two differentiable heads serve complementary roles. Attention--linear
provides an efficient neural reference and strong raw probability
calibration. ODST provides a structured decision layer that exposes feature
selection over the learned representation, route probabilities, reached
leaves, individual tree outputs, and tree-level contributions.

\smallskip
Reference-centred ODST contributions reconstructed the model readout
exactly. Top-tree ablation exceeded matched random-tree ablation across all
three final CERT r5.2 seeds at \(k=3\), \(5\), and \(10\). This supports the
faithfulness of the reference-centred tree-contribution rankings under the
defined intervention; routes and reached leaves remain traceable internal
outputs rather than independently validated attribution levels. Ranking by
raw absolute tree output, together with centred ranking against the median
training reference, is retained separately as a sensitivity result, while
the disclosed post-hoc test analysis provides corroboration rather than
primary evidence.

\smallskip
The explanation procedure also transferred from CERT r4.2 to r5.2.
Reference-centred faithfulness remained consistent, while the most
influential original features and source channels varied by dataset. This
is an appropriate portability result: the explanation mechanism transferred
without assuming that different CERT versions must contain identical
behavioural patterns.

\smallskip
Grouped out-of-fold temperature scaling substantially improved ODST
probability quality. For seed 42, the Brier score improved from 0.175 to
0.004 and expected calibration error from 0.413 to 0.005. Ranking remained
exact within each fold, although fold-specific temperatures limit direct
comparison of globally combined out-of-fold ranks.

\smallskip
Alert consolidation provided an operational view of the overlapping
sequence outputs. The 661 seed-42 sequence alerts represented 36 users and
42 one-day-gap analytical episodes. ODST and attention--linear remained
close at the predefined alert budgets, showing that the predictive
near-parity also extended to the evaluated alert-volume conditions.
Episodes are interpreted as temporal analytical groups rather than as
independently verified incidents.

\subsection{Computational and Degraded-Input Evidence}
\label{subsec:discussion_efficiency}
The component audit localised the main inference cost to the ODST head,
which accounted for approximately 72--73\% of device-resident latency. The
eight-tree study halved the head parameter count while preserving PR-AUC,
F1-score, and reference-centred faithfulness. Its batch-32 latency was
nevertheless 7.4\% higher than that of the protected 16-tree model. The
smaller head therefore demonstrated removable parameter capacity but did
not satisfy the predefined runtime objective.

\smallskip
This result gives a clear engineering direction. Further efficiency work
should focus on ODST execution, including sparse selection, routing, leaf
evaluation, and kernel organisation, rather than assuming that reducing
tree count will automatically reduce latency.

\smallskip
Controlled source ablation and Gaussian noise provided a separate account
of model dependency. Source effects were measurable across CERT versions,
although their ordering varied with the dataset. Continuous noise produced
dose-dependent PR-AUC degradation, while route and leaf assignments remained
stable under low noise for cases whose predicted class was unchanged.

\smallskip
These experiments establish reproducible targets for data-quality controls
and future hardening. They are controlled degraded-input studies rather
than evaluations of adversarial evasion or training-time poisoning.

\subsection{Overall Contribution and Scope}
\label{subsec:discussion_contribution}
The study makes four connected contributions, matching those introduced in
Section~\ref{sec:introduction}. It establishes the predictive value of
chronology and accumulated behavioural history; demonstrates stable
teacher-anchored end-to-end optimisation, reproducible across three
predefined r5.2 seeds; compares sequence and flattened models using equal
information; and provides an intervention-based evaluation of
reference-centred ODST tree-contribution faithfulness, complemented by
traceable route and leaf outputs, probability-calibration analysis, alert
consolidation, component-level latency profiling, and controlled
degraded-input analysis.

\smallskip
The resulting component-level interpretation is clear. Bi-LSTM--attention
provides the principal temporal predictive contribution. Teacher anchoring
provides the optimisation and within-dataset reproducibility contribution.
ODST serves as an explanation-oriented differentiable head.
Attention--linear remains the efficiency- and raw-calibration-oriented
neural reference, while RF and XGBoost remain demanding engineered-tabular
comparators.

\smallskip
The evidence is based on synthetic CERT benchmark activity rather than live
organisational deployment. The bidirectional encoder supports completed
behavioural windows, and overlapping sequence alerts are not independent
incidents. Human-analyst usefulness, live deployment, adversarial evasion,
and training-time poisoning were not evaluated.

\smallskip
The next bounded study is a CERT r6.2 metadata, ground-truth,
feature-interface, and temporal-partition feasibility audit without model
training or inference. A later extreme-rarity experiment should proceed only
after a defensible partition has been reviewed and frozen. This preserves
the completed evidence while extending the research through a controlled
cross-version study.